% =========================================================
%  NSFC Style Configuration (风格配置文件)
%  说明：此文件仅用于控制排版风格开关，核心逻辑在 config.tex
% =========================================================

% ---------------------------------------------------------
%  开关 1：正文字体设置
% ---------------------------------------------------------
% true  = 正文强制使用【楷体】
% false = 正文使用标准【宋体】
\newif\ifUseKaitiBody
\UseKaitiBodyfalse  % <--- 修改这里：true 或 false

% ---------------------------------------------------------
%  开关 2：标题字体选择
% ---------------------------------------------------------
% 定义各级标题（一级、二级）使用的字体
% 常用选项：\kaishu (楷体), \heiti (黑体), \songti (宋体)
\newcommand{\HeadingFont}{\mysong} 

% ---------------------------------------------------------
%  开关 3：标题编号风格
% ---------------------------------------------------------
% 请输入数字选择一种风格：
% 1 = 标准简洁: 1  / 1.1 / 1.1.1 
% 2 = Word风格: 1. / 1.1 / 1.1.1 (一级标题带个点)
% 3 = 中文风格: (一) / 1. / (1)  (公文风)
\newcommand{\NumberingStyle}{2} % <--- 修改这里：1, 2, 或 3

% ---------------------------------------------------------
%  开关 4：交叉引用与链接颜色
% ---------------------------------------------------------
% 定义参考文献标号、目录、公式引用的颜色。
% 打印版（黑白纸质）强烈建议用 black。电子版可用 blue, red, green, darkgray 等
\newcommand{\ThemeColor}{blue} % <--- 修改这里：输入你想要的颜色单词


% ---------------------------------------------------------
%  开关 5：全局布局与间距调控
% ---------------------------------------------------------
% 设置说明：
% 1. 段间距 (UserParskip): 建议值 0.3ex 到 0.8ex。若篇幅吃紧可设为 0pt。
\newcommand{\UserParskip}{0.5ex} % <--- 修改这里：段间距数值 (如 0.5ex, 3pt)

% 2. 列表模式 (UserListMode): 
%    0 = 极致紧凑: 列表项内外完全不留白 (nosep)，最省空间。
%    1 = 宽松美观: 列表间距随段间距自动按比例缩放，呼吸感强。
\newcommand{\UserListMode}{0}    % <--- 修改这里：0 或 1

%-----------------------------------------------
% 自定义命令（自定义颜色加粗）
%-----------------------------------------------
\definecolor{myemphasize}{RGB}{0, 0, 64} % 经典学术蓝，比纯黑更有质感#000640
% 定义自定义加粗命令 \mybf
\newcommand{\mybf}[1]{{\color{myemphasize}\textbf{#1}}}
